% Generated from locale/tr/content/methods/proofs/proof-by-contradiction.tex; do not hand-edit.


\olfileid[tr]{mth}{prf}{con}

\olsection{Çelişki Yoluyla Kanıt}

İlk olarak çelişki yoluyla kanıt, olumsuz iddiaları kanıtlamak için
kullanılan bir çıkarım kalıbıdır. Bir $p$ iddiasının \emph{yanlış}
olduğunu, yani~$\lnot p$'yi göstermek istediğinizi düşünün. En umut
verici strateji (a) $p$'nin doğru olduğunu varsaymak ve (b) bu
varsayımın yanlış olduğunu bildiğiniz bir şeye götürdüğünü göstermektir.
``Yanlış olduğu bilinen şey'', $p$'nin kendisiyle ya da ele aldığınız
genel iddianın başka bir hipoteziyle çatışan---çelişen---bir sonuç
  olabilir. Örneğin ``$q$ ise $\lnot p$''nin bir kanıtı, $q$'nun doğru
  olduğunu varsayıp buradan~$\lnot p$'yi kanıtlamayı içerir. $\lnot p$'yi
  çelişki yoluyla kanıtlarsanız, bu $p$'yi $q$'ya ek olarak varsaymak
  demektir. $\lnot q$'yu $p$'den kanıtlayabilirseniz, $p$ varsayımının
diğer varsayımınız~$q$ ile çelişen bir şeye götürdüğünü göstermiş
olursunuz; çünkü $q$ ve $\lnot q$ birlikte doğru olamaz. Elbette
çelişkiyi kanıtlarken tanımları açmanın yanı sıra başka çıkarım
kalıplarını da kullanmanız gerekir. Bir örnek ele alalım.

\begin{prop}
  $A \subseteq B$ ve $B = \emptyset$ ise $A$ için !!{element}s
  bulunmaz.
\end{prop}

\begin{proof}
  $A \subseteq B$ ve $B = \emptyset$ olduğunu varsayalım. $A$ için
  !!{element}s bulunmadığını göstermek istiyoruz.
  \begin{quote}
    Bu koşullu bir iddia olduğundan önbileşeni varsayar ve artbileşeni
    kanıtlamak isteriz. Artbileşen şudur: $A$ için !!{element}s
    bulunmaz. Bunu biraz daha açık hâle getirebiliriz: bir~$x \in A$
    bulunması söz konusu değildir.
  \end{quote}
  $A$ için !!{element}s bulunmaması, ancak ve ancak bir~$x$ için $x \in A$
  olması söz konusu değilse geçerlidir.
  \begin{quote}
    Böylece kanıtlamak istediğimiz şeyin aslında olumsuz bir iddia
    $\lnot p$, yani bir $x \in A$ bulunmasının söz konusu olmadığı
    olduğunu belirledik. Çelişki yoluyla kanıtı kullanmak için buna
    karşılık gelen olumlu iddiayı~$p$'yi, yani bir $x \in A$ bulunduğunu
    varsaymalı ve buradan bir çelişki kanıtlamalıyız. Çelişki yoluyla
    kanıt yaptığımızı, ``tersini varsayalım'' ya da yalnızca ``öyle
    olmadığını varsayalım'' yazıp ardından~$p$ varsayımını belirterek
    gösteririz.
  \end{quote}
  Öyle olmadığını varsayalım: bir $x \in A$ vardır.
  \begin{quote}
    Artık bir çelişki elde etmek için kullanacağımız yeni varsayım
    budur. İki varsayımımız daha var: $A \subseteq B$ ve $B =
    \emptyset$. Birincisi bize $x \in B$ sonucunu verir:
  \end{quote}
  $A \subseteq B$ olduğundan $x \in B$.
  \begin{quote}
    Fakat $B = \emptyset$ olduğundan, $B$'deki her !!{element} (örneğin
    $x$) aynı zamanda $\emptyset$ içinde bulunan bir !!a{element} olmalıdır.
  \end{quote}
  $B = \emptyset$ olduğundan $x \in \emptyset$. Bu bir çelişkidir;
  çünkü tanım gereği $\emptyset$ için !!{element}s bulunmaz.
  \begin{quote}
    Bu, kanıtı zaten tamamlar: kurduğumuz varsayımlardan ihtiyacımız
    olan şeye (bir çelişkiye) ulaştık; demek ki bu varsayımların hepsi
    birden doğru olamaz. İlk iki varsayım ($A \subseteq B$ ve $B =
    \emptyset$) tartışmalı olmadığına göre, yanlış olması gereken son
    ortaya konan varsayımdır (bir $x \in A$ vardır). Ama eksiksiz olmak
    istersek bunu açıkça yazabiliriz.
  \end{quote}
  Öyleyse bir $x \in A$ bulunduğu varsayımımız yanlış olmalıdır;
  dolayısıyla çelişki yoluyla kanıt gereği $A$ için !!{element}s
  bulunmaz.
\end{proof}

Her olumlu iddia önemsiz biçimde bir olumsuz iddiaya denktir: $p$, ancak
ve ancak $\lnot\lnot p$ ise geçerlidir. Dolayısıyla çelişki yoluyla
kanıtlar, olumlu iddiaları ``dolaylı olarak'' belirlemek için de şöyle
kullanılabilir: $p$'yi kanıtlamak için onu $\lnot\lnot p$ olumsuz
iddiası olarak okuyun. $\lnot p$'den bir çelişki kanıtlayabilirsek,
çelişki yoluyla kanıtla $\lnot\lnot p$'yi, dolayısıyla~$p$'yi belirlemiş
oluruz.

Son örnekte olumsuz bir iddiayı, yani $A$ için !!{element}s
bulunmadığını kanıtlamayı hedefledik. Bu nedenle çelişki yoluyla kanıt
amacıyla yaptığımız varsayım (yani bir $x \in A$ bulunduğu), olumlu bir
iddiaydı. Bize üzerinde çalışabileceğimiz bir şey, yani hakkında akıl
yürütmeyi sürdürdüğümüz varsayımsal $x \in A$'yı verdi; sonunda $x \in
\emptyset$ sonucuna ulaştık.

Olumlu bir iddiayı dolaylı olarak kanıtlarken çelişki yoluyla kanıt
amacıyla yapacağınız varsayım olumsuz olur. Fakat çoğu zaman olumlu bir
  iddiayı olumsuz bir iddia, olumsuz bir iddiayı da olumlu bir iddia
  olarak kolayca yeniden ifade edebilirsiniz. Bir önceki kanıtımız,
  ``$A = \emptyset$''i ``$A$ için !!{element}s bulunmaz'' olumsuz
  artbileşeni yerine kanıtlamış olsaydık özünde aynı olurdu. ($=$ tanımı gereği
``$A = \emptyset$'' genel bir iddiadır; çünkü açıldığında ``$A$'da
bulunan her !!{element}, aynı zamanda $\emptyset$'te bulunan bir
!!{element} olur ve bunun tersi de geçerlidir'' olur.) Ama bunun,
``değil: bir $x \in A$ vardır'' olumsuz
iddiasına denk olduğu kolayca görülür.

Dolayısıyla bazen $\lnot p$ ile varsayım olarak çalışmak, $p$'yi
doğrudan kanıtlamaktan daha kolaydır. Doğrudan kanıt aynı derecede basit
hatta daha basit olsa bile (sonraki örneklerde olduğu gibi), bazı
insanlar dolaylı ilerlemeyi yeğler. Çifte olumsuzlama kafanızı
karıştırıyorsa, bir iddianın çelişki yoluyla kanıtını, \emph{karşıt}
iddiadan bir çelişkinin kanıtı olarak düşünün. Buna göre $\lnot p$'nin
çelişki yoluyla kanıtı, $p$ varsayımından bir çelişkinin kanıtıdır;
$p$'nin çelişki yoluyla kanıtı ise $\lnot p$'den bir çelişkinin
kanıtıdır.

\begin{prop}
$A \subseteq A \cup B$.
\end{prop}

\begin{proof}
  $A \subseteq A \cup B$ olduğunu göstermek istiyoruz.
  \begin{quote}
    Görünüşte bu olumlu bir iddiadır: her $x \in A$, aynı zamanda $A
    \cup B$'dedir. Bunun olumsuzlaması şudur: bazı $x \in A$ öğeleri
    $\notin A \cup B$'dir. Böylece bu olumsuzlanmış iddiayı varsayıp onun
    bir çelişkiye götürdüğünü göstererek iddiayı dolaylı biçimde
    kanıtlayabiliriz.
    \end{quote}
  Öyle olmadığını, yani $A \nsubseteq A \cup B$ olduğunu varsayalım.
  \begin{quote}
    $A \subseteq A \cup B$'nin bir tanımı vardır: her $x \in A$, aynı
    zamanda $\in A \cup B$'dir. $A \nsubseteq A \cup B$'nin ne anlama
    geldiğini anlamak için açılmış tanım üzerinde bazı temel mantıksal
    işlemler yapmalıyız: her $x \in A$'nın aynı zamanda $\in A \cup B$
    olması yanlışsa, ancak ve ancak \emph{bazı}~$x \in A$ öğeleri
    $\notin A \cup B$ ise durum böyledir. (Bu, çok dikkatli olmanız gereken bir
    yerdir: öğrencilerin birçok çelişki yoluyla kanıt girişimi, ``bütün
    $A$'lar $B$'dir'' gibi bir iddianın olumsuzlamasını yanlış çözümlediği
    için başarısız olur.) Başka bir deyişle $A \nsubseteq A \cup B$,
    ancak ve ancak bir $x$ bulunup $x \in A$ ve $x \notin A \cup B$ ise
    geçerlidir. Bundan sonrası kolaydır.
  \end{quote}
  Öyleyse bir $x \in A$ vardır ve $x \notin A \cup B$. $\cup$ tanımı
  gereği $x \in A \cup B$, ancak ve ancak $x \in
  A$ ya da $x \in B$ ise geçerlidir. $x \in A$ olduğundan $x \in A
  \cup B$ elde ederiz. Bu, $x \notin A \cup B$ varsayımıyla çelişir.
\end{proof}

\begin{prob}
\emph{Dolaylı olarak} $A \cap B \subseteq A$ olduğunu kanıtlayın.
\end{prob}

\begin{prop}
$A \subseteq B$ ve $B \subseteq C$ ise $A \subseteq C$.
\end{prop}

\begin{proof}
  $A \subseteq B$ ve $B \subseteq C$ olduğunu varsayalım. $A
  \subseteq C$ olduğunu göstermek istiyoruz.
  \begin{quote}
    Dolaylı ilerleyelim: kanıtlamak istediğimiz şeyin olumsuzlamasını
    varsayalım.
  \end{quote}
  Öyle olmadığını, yani $A \nsubseteq C$ olduğunu varsayalım.
  \begin{quote}
    Daha önce olduğu gibi $A \nsubseteq C$'nin, her $x \in A$'nın aynı
    zamanda $\in C$ olmamasıyla, yani bazı $x \in A$ öğelerinin $\notin C$
    olmasıyla denk olduğunu görürüz. Kaygılanmayın; biraz alıştırmayla
    böyle olumsuzlamaları açmak için artık çok düşünmeniz gerekmeyecek.
  \end{quote}
  Başka bir deyişle bir~$x$ vardır; $x \in A$ ve $x \notin C$.
  \begin{quote}
    Artık çelişkimize ulaşmak için bunu kullanabiliriz. Elbette bunu
    yapmak üzere öteki iki varsayımı da kullanmamız gerekecek.
  \end{quote}
  $A \subseteq B$ olduğundan $x \in B$. $B \subseteq C$ olduğundan $x
  \in C$. Fakat bu $x \notin C$ ile çelişir.
\end{proof}

\begin{prop}
$A \cup B = A \cap B$ ise $A = B$.
\end{prop}

\begin{proof}
  $A \cup B = A \cap B$ olduğunu varsayalım. $A = B$ olduğunu göstermek
  istiyoruz.
  \begin{quote}
    Başlangıç artık bildik bir yoldur:
  \end{quote}
  Çelişkiye ulaşmak üzere $A \neq B$ olduğunu varsayalım.
  \begin{quote}
    Çelişki yoluyla kanıt için varsayımımız $A \neq B$'dir. $A = B$,
    ancak ve ancak $A \subseteq B$ ve $B \subseteq A$ ise geçerli
    olduğundan, $A \neq B$'nin ancak ve ancak $A \nsubseteq B$ \emph{ya
    da} $B \nsubseteq A$ ise geçerli olduğu sonucuna varırız.
    (Olumsuzlamaları işlerken dikkatli olmanın ne kadar önemli olduğuna
    bakın!{}) Bu ayrışımdan bir çelişki kanıtlamak için durumlara göre
    kanıt kullanır ve her durumda bir çelişkinin çıktığını gösteririz.
  \end{quote}
  $A \neq B$, ancak ve ancak $A \nsubseteq B$ ya da $B \nsubseteq A$
  ise geçerlidir. Durumları ayıralım.
  \begin{quote}
    İlk durumda $A \nsubseteq B$'yi, yani bazı $x$ için $x \in A$ fakat
    $\notin B$ olduğunu varsayarız. $A \cap B$, $A$ ile $B$'nin ortak
    !!{element}s kümesi olarak tanımlandığından bir şey bunlardan birinde
    değilse kesişimde de değildir. $A \cup B$, $A$ ile birlikte $B$'dir;
    dolayısıyla ikisinden herhangi birinde bulunan her şey birleşimde
    de bulunur. Bu bize $x \in A \cup B$ fakat $x \notin A \cap B$
    olduğunu ve dolayısıyla $A \cap B \neq A \cup B$ olduğunu söyler.
  \end{quote}

  Durum 1: $A \nsubseteq B$. O zaman bazı $x$ için $x \in A$ fakat $x
  \notin B$. $x \notin B$ olduğundan $x \notin A \cap B$. $x \in A$
  olduğundan $x \in A \cup B$. Öyleyse $A \cap B \neq A \cup B$;
  bu, $A \cap B = A \cup B$ varsayımıyla çelişir.

  Durum 2: $B \nsubseteq A$. O zaman bazı $y$ için $y \in B$ fakat $y
  \notin A$. Daha önce olduğu gibi $y \in A \cup B$ fakat $y \notin A
  \cap B$ elde ederiz; dolayısıyla $A \cap B \neq A \cup B$, bu da $A
  \cap B = A \cup B$ ile çelişir.
\end{proof}

